\chapter{Decisiones, juegos y programación dinámica}
\label{chap:decisiones}

\section{Elegir bajo incertidumbre}
Una decisión combina alternativas controlables con estados del entorno que el
analista no controla. Una tabla de pagos organiza el beneficio $a_{ij}$ de la
alternativa $i$ en el estado $j$. Cuando se conocen probabilidades previas $p_j$,
el criterio de valor esperado selecciona
\[
 i^*\in\arg\max_i\sum_j p_j a_{ij}.
\]
Este criterio supone pagos comparables y neutralidad frente al riesgo. Una
solución que maximiza el beneficio medio puede incluir pérdidas importantes en
un estado concreto. QSB no transforma automáticamente dinero en utilidad ni
incorpora la aversión al riesgo del decisor.

\subsection{Tabla de pagos e información}
El archivo \path{dec_payoff.json} plantea un lanzamiento con demanda Alta
(probabilidad $0{,}4$) o Baja ($0{,}6$). Los pagos son:
\begin{center}
\begin{tabular}{lrr}
\toprule
Decisión & Alta & Baja\\
\midrule
Lanzar & 100 & $-40$\\
Esperar & 20 & 10\\
\bottomrule
\end{tabular}
\end{center}
Sin estudiar el mercado, Lanzar obtiene $0{,}4(100)+0{,}6(-40)=16$, frente a
14 de Esperar. Con información perfecta se elegiría Lanzar en Alta y Esperar
en Baja: $0{,}4(100)+0{,}6(10)=46$. El valor esperado de información perfecta
(EVPI) es $46-16=30$, un techo para el precio de información bajo este modelo.

El estudio produce un indicador Positivo o Negativo. La matriz de
\texttt{indicatorLikelihoods} usa filas de indicadores y columnas de estados:
$P(+\mid Alta)=0{,}8$ y $P(+\mid Baja)=0{,}2$. Cada columna representa una
distribución sobre los posibles indicadores. No son probabilidades posteriores.
La política con estudio es Lanzar si sale Positivo y Esperar si sale Negativo.
Su valor esperado es $33{,}6$; por tanto, EVSI (valor de información muestral)
es $33{,}6-16=17{,}6$. Este importe no descuenta el coste del estudio. Si el
estudio costase 20, no convendría comprarlo con el criterio de valor esperado.

\begin{lstlisting}
swift run qsb validate-decision-json manual/examples/dec_payoff.json
swift run qsb solve-decision-json manual/examples/dec_payoff.json --backend native
\end{lstlisting}
La salida contiene \texttt{priorExpectedValues}, \texttt{indicatorAnalyses},
\texttt{bestPriorDecision}, EVSI y EVPI. Las claves conservan sus nombres en
inglés. El solver exige probabilidad mayor que $10^{-12}$ para cada indicador;
si no se cumple, falla la resolución en vez de inventar una posterior.

\subsection{Actualización bayesiana}
La regla de Bayes convierte probabilidad previa y verosimilitud en probabilidad
posterior:
\[
 P(s_j\mid e_k)=\frac{P(e_k\mid s_j)p_j}{\sum_\ell P(e_k\mid s_\ell)p_\ell}.
\]
El archivo \path{dec_bayes.json} separa esta operación de cualquier tabla de
pagos. Usa los mismos estados y el mismo estudio. QSB obtiene
$P(+)=0{,}44$, $P(-)=0{,}56$, y las posteriores:
\begin{center}
\begin{tabular}{lrr}
\toprule
Evidencia & Alta & Baja\\
\midrule
Positivo & $0{,}727273$ & $0{,}272727$\\
Negativo & $0{,}142857$ & $0{,}857143$\\
\bottomrule
\end{tabular}
\end{center}
\begin{lstlisting}
swift run qsb validate-decision-json manual/examples/dec_bayes.json
swift run qsb solve-decision-json manual/examples/dec_bayes.json
\end{lstlisting}
\texttt{likelihoods} vuelve a usar filas de resultados y columnas de estados.
Cambiar el orden de los nombres exige cambiar también las filas o columnas
correspondientes. Las probabilidades deben ser finitas, no negativas y
consistentes con las distribuciones representadas. Bayes actualiza creencias;
por sí solo no indica qué acción tomar.

\subsection{Árboles y análisis hacia atrás}
Un árbol hace explícito el orden temporal de decisiones y observaciones.
\texttt{rootID} identifica la raíz, cada nodo tiene un \texttt{id} único y sus
\texttt{childIDs} indican sucesores. Los tipos son \texttt{decision},
\texttt{chance} y \texttt{terminal}. Las probabilidades de las ramas de azar
se guardan en los nodos hijos, no en una lista del padre.

El análisis hacia atrás asigna al terminal su pago. Un nodo de decisión suma
su pago local al máximo valor de sus hijos; uno de azar suma su pago local al
valor esperado de los hijos. Los pagos locales se acumulan a lo largo de la
trayectoria: no repita en las hojas un coste que ya introdujo en un antecesor.

En \path{dec_tree.json}, Invertir cuesta 5, da 40 con probabilidad $0{,}6$
y $-10$ con probabilidad $0{,}4$. Su valor es
$-5+0{,}6(40)+0{,}4(-10)=15$; Conservar da 12. El valor de la raíz es 15 y
la política selecciona Invertir.
\begin{lstlisting}
swift run qsb validate-decision-json manual/examples/dec_tree.json
swift run qsb solve-decision-json manual/examples/dec_tree.json
\end{lstlisting}
\begin{figure}[H]
\centering
\begin{tikzpicture}
\node[rectangle,draw,minimum size=6mm,label=above:{\small Elegir}] (R) at (0,0) {};
\node[qsb node,label=above:{\small Invertir},label=below:{$-5$}] (I) at (4,1) {};
\node[qsb terminal] (K) at (4,-1.5) {Conservar\\12};
\node[qsb terminal] (E) at (8,2) {Exito\\40};
\node[qsb terminal] (F) at (8,-0.2) {Fracaso\\$-10$};
\draw[qsb selected] (R) -- (I);
\draw[qsb alternative] (R) -- (K);
\draw[qsb arc] (I) -- node[qsb label,above] {$0.6$} (E);
\draw[qsb arc] (I) -- node[qsb label,below] {$0.4$} (F);
\end{tikzpicture}

\caption{Árbol de decisión: el cuadrado permite elegir y el círculo representa azar. La rama gruesa es la política elegida. Las hojas muestran pagos; el costo local de invertir se añade una sola vez. Los nombres conservan los identificadores JSON.}
\label{fig:decision-tree}
\end{figure}

Revise \texttt{nodeValues} para cada nodo y \texttt{policy} para las decisiones.
El método maximiza pagos: para minimizar costes puede usar sus negativos,
explicando la convención. QSB advierte y normaliza probabilidades de hijos de
azar cuya suma positiva difiere de uno, por compatibilidad con entradas
redondeadas. Esa reparación no sustituye una revisión de los datos. Evite
ciclos, identificadores inexistentes y probabilidades incoherentes.

\section{Juegos de suma cero}
En un juego de suma cero, el pago del jugador de filas es la pérdida del jugador
de columnas. El primero maximiza y el segundo minimiza. Las estrategias mixtas
son distribuciones de probabilidad sobre acciones; no representan cantidades
fraccionarias de una acción en una partida.

Para una matriz $A$, QSB resuelve problemas lineales equivalentes. El jugador
de filas busca $\max v$ sujeto a $A^Tp\geq v\mathbf{1}$,
$\sum_i p_i=1$, $p_i\geq0$. El jugador de columnas busca el problema dual de
minimización. Esto presupone elección simultánea y pagos conocidos.
\path{dec_game.json} usa
\[
 A=\begin{pmatrix}1&-1\\-1&1\end{pmatrix}.
\]
Ambos jugadores asignan probabilidad $0{,}5$ a A y B, y el valor del juego es
cero. La salida numérica puede mostrar $4{,}44\times10^{-16}$ por redondeo.
\begin{lstlisting}
swift run qsb validate-decision-json manual/examples/dec_game.json
swift run qsb solve-decision-json manual/examples/dec_game.json
\end{lstlisting}
Lea \texttt{rowStrategy}, \texttt{columnStrategy} y \texttt{value}. El método
nativo usa el solver lineal compartido. No interpreta esta solución como un
equilibrio de un juego general de suma no nula ni de un juego secuencial.

\section{Programación dinámica: estado, etapa y acción}
La programación dinámica divide un problema en etapas. El estado resume la
información necesaria para valorar las decisiones futuras; la acción cambia
ese estado y produce un coste o retorno local. El principio de optimalidad
permite construir una política óptima a partir de subproblemas. Un estado que
omite información relevante rompe ese principio.

QSB implementa tres modelos discretos y deterministas: mochila acotada,
ruta acíclica por etapas y producción con inventario. Las entradas usan
\texttt{kind} y \texttt{model}. Las salidas incluyen un resultado tipado,
\texttt{assumptions} y \texttt{trace}. Cada fila de la traza describe etapa,
estado, acción, siguiente estado y valor \emph{local}; ese valor no es
necesariamente la función de valor óptimo del subproblema.

\subsection{Mochila acotada}
Cada artículo $i$ consume $w_i$ unidades de capacidad y aporta $r_i$ por unidad;
hay como máximo $b_i$ unidades disponibles. La formulación es
\[
 \max\sum_i r_i x_i,\qquad \sum_i w_i x_i\leq C,
 \qquad x_i\in\{0,\ldots,b_i\}.
\]
Una recurrencia con capacidad restante $c$ es
$F_i(c)=\max_q\{q r_i+F_{i+1}(c-qw_i)\}$, para cantidades disponibles que
caben. La capacidad y los consumos son enteros; los retornos pueden ser reales
finitos. No convierta medidas físicas continuas a enteros sin justificar la
unidad de discretización.

En \path{dec_knapsack.json}, $C=7$ y los artículos
$(b_i,w_i,r_i)$ son A $(2,3,5)$, B $(3,2,3)$ y C $(1,4,7)$.
QSB selecciona una A y una C: capacidad 7 y retorno 12. La alternativa una A
y dos B ocupa también 7, pero retorna 11. La selección codifica cantidades;
la traza incluye también la decisión de tomar cero unidades de B.
\begin{lstlisting}
swift run qsb validate-dp-json manual/examples/dec_knapsack.json
swift run qsb solve-dp-json manual/examples/dec_knapsack.json
\end{lstlisting}
En macOS, Bounded Knapsack permite editar capacidad, disponibilidades,
consumos y retornos, así como copiar o pegar la tabla. Summary ofrece métricas,
barras y una tabla de selección; la traza permite explicar la política.

\subsection{Ruta por etapas (Stagecoach)}
La lista \texttt{nodes} declara un orden topológico: el primer nodo es origen
y el último destino. Todo arco debe avanzar en ese orden y tener coste finito
no negativo. El método aplica
\[
 V(i)=\min_{(i,j)}\{c_{ij}+V(j)\},\qquad V(\text{destino})=0.
\]
\path{dec_route.json} contiene Origen--Norte (4), Origen--Sur (2),
Norte--Destino (3) y Sur--Destino (6). El resultado es
Origen--Norte--Destino, con coste 7; elegir el arco inicial más barato daría 8.
\begin{lstlisting}
swift run qsb validate-dp-json manual/examples/dec_route.json
swift run qsb solve-dp-json manual/examples/dec_route.json
\end{lstlisting}
Un grafo estructuralmente válido puede carecer de ruta hasta el destino; eso
se detecta al resolver. Este modelo no acepta ciclos ni arcos hacia atrás:
para otros grafos consulte el capítulo de redes.

\subsection{Producción e inventario}
El estado al inicio de un periodo es el inventario $I_{t-1}$; la acción es
producción entera $P_t$. Se satisface demanda determinista $d_t$ con
$I_t=I_{t-1}+P_t-d_t\geq0$. Las capacidades limitan $P_t$ e $I_t$.
QSB supone inventario inicial y final cero, sin pedidos pendientes. Minimiza
\[
 \sum_t\left[K_t\mathbf{1}_{P_t>0}+c_tP_t+h_tI_t\right].
\]
El coste de posesión se aplica al inventario final del periodo, no al inicial.
La recurrencia evalúa producciones factibles y añade el coste óptimo futuro.

\path{dec_production.json} tiene demandas 2 y 3, capacidad de producción 5,
capacidad de almacén 3, preparación 10, coste unitario 2 y posesión 1 en ambos
periodos. QSB produce 5 al inicio, conserva 3 y no produce en el segundo
periodo. El coste es $10+2(5)+1(3)=23$, frente a 30 si prepara y produce en
cada periodo. Las decisiones detallan inventario inicial, producción,
demanda, inventario final y coste local.
\begin{lstlisting}
swift run qsb validate-dp-json manual/examples/dec_production.json
swift run qsb solve-dp-json manual/examples/dec_production.json
\end{lstlisting}
La validación de campos no prueba que exista un plan factible. Capacidades
insuficientes pueden causar fallo de resolución. Este modelo no representa
varios productos, inventario inicial configurable, faltantes o demanda aleatoria.

\section{Uso, metadatos y límites comunes}
Las cuatro variantes de decisiones tienen editores nativos en macOS. El árbol
permite editar nodos y enlaces; eliminar hojas o añadir terminales conserva una
estructura más fácil de revisar. Stagecoach y producción/inventario pueden
editarse mediante el formulario estructurado o JSON; el resultado tiene vista
nativa de ruta o series por periodo. Use Validation antes de Run y consulte el
JSON cuando necesite conservar todos los campos.

El backend nativo de decisiones declara \texttt{exact}: calcula el criterio
especificado con aritmética numérica, dentro del modelo y la escala educativa.
Las tres variantes dinámicas declaran \texttt{fixtureScale}; sus políticas son
las obtenidas por las recurrencias implementadas, con discretización y consumo
de memoria que crecen con capacidad, inventario y número de etapas. No existe
una promesa de rendimiento industrial. \texttt{--backend validate} devuelve
validación y no una solución; el backend externo no está disponible para estas
familias. Guarde el modelo junto a la salida de programación dinámica: su
documento de solución conserva título y supuestos, pero no el modelo completo.

\section{Ejercicios y comprobaciones}
\begin{enumerate}
 \item Cambie $P(Alta)$ de $0{,}4$ a $0{,}3$, ajustando también Baja.
 Calcule primero la mejor decisión sin información y vuelva a resolver.
 \textquestiondown{}Por qué puede cambiar EVSI aunque no cambie la calidad del estudio?
 \item Calcule el precio máximo del estudio del ejemplo y compárelo con EVPI.
 Explique qué supuestos impedirían usar ese precio en una decisión real.
 \item En el árbol, eleve el coste de Invertir a 9. El valor de esta rama será
 11 y la decisión de raíz pasará a Conservar, con valor 12.
 \item En la mochila, cambie la capacidad a 6. Verifique el retorno 10 con
 dos unidades A y explique por qué el cociente retorno/consumo no basta para
 asegurar optimalidad.
 \item En la ruta, haga que Sur--Destino cueste 4. La ruta por Sur costará 6.
 Añada después un arco hacia atrás y examine el diagnóstico de validación.
 \item Reduzca el coste de preparación de producción a 1 en ambos periodos.
 Producir 2 y 3 por separado cuesta 12, frente a 14 al acumular inventario.
 Compare cada balance de inventario y la traza.
\end{enumerate}
